\documentclass[11pt]{article}
\usepackage[margin=1in]{geometry}
\usepackage{amsmath,amssymb,amsthm}
\usepackage{booktabs}
\usepackage{graphicx}
\usepackage{hyperref}
\usepackage{xcolor}
\newtheorem{proposition}{Proposition}
\newtheorem{remark}{Remark}
\title{\bf Gradient-Weighted Compressed-Response Screening for\\
Branch-and-Bound in Single-Track Train Timetabling}
\author{Transportation+AI Lab \\ \small (working draft --- single-track results)}
\date{\today}

\begin{document}
\maketitle

\begin{abstract}
Exact branch-and-bound (B\&B) and branch-and-price for train timetabling spend most of
their effort \emph{evaluating} branching candidates: at each node, strong branching may
solve many child relaxations. We study whether the low-rank structure of the
\emph{solver's response field} can decide which candidates deserve that exact effort. We
build a faithful single-track timetabling B\&B (meet--pass precedence branching with a
crossing-conflict enhanced lower bound, after Zhou and Zhong 2007) and treat it as a
data-generating process: every node emits candidate meet--pass actions and their
primal--dual responses. We derive a \emph{compressed-response screen}: branch candidates are
mapped to low-rank response coordinates and scored, the top of a \emph{statistical tie
cluster} $K'$ is certified exactly, and the certified decision is preserved. The central
theoretical point is that the screening score must be \emph{gradient-weighted}
($\sum_i \sigma_i (u_i^\top\nabla S)(v_i^\top b_a)$): a discarded response direction harms
the ranking only if the objective gradient points along it. Empirically, on single-track
instances: (i) the crossing-conflict lower bound reduces B\&B nodes by $7.5\%$ to $35.6\%$
as trains grow from $6$ to $14$ (exact optimum preserved); (ii) a learned, gradient-aware
screen contains the full-strong-branching action in its top-$5$ on $92\%$ of nodes and
\emph{transfers} to unseen instances, clearly beating an unsupervised spectral score and
pseudo-cost; (iii) the per-node response tensor is low multilinear rank (the metric mode is
rank one). We are explicit about scope: on single-track, candidate sets are small and the
cheap earliest-conflict rule already minimizes the tree, so the payoff is screening fidelity
and the lower-bound node reduction---not a per-solve speedup. The speedup regime (large
candidate sets, true LP duals, route/track choice) is the $N$-track space-time network, for
which we have a working instance and which we leave as the certified next step.
\end{abstract}

\section{Introduction}
Railway timetabling selects, for each train, a conflict-free trajectory over shared
time--space resources (track segments, sidings, junctions) subject to running times, dwell,
meet--pass order, and capacity. Exact methods---branch-and-bound (B\&B) and branch-and-price
(B\&P)---are natural, but their cost is dominated by \emph{branching}: at a search node the
algorithm faces many heterogeneous candidate decisions, and full strong branching evaluates
each by (partially) solving its child relaxations. Pseudo-cost and rule-based branching are
cheap but ignore node-specific conflict structure; learned (GNN) branching imitates strong
branching but is a black box and offers no certificate.

This paper asks a structural question: \emph{the solver itself generates data---primal
trajectories, conflict structure, child-bound improvements---so what is the low-rank
structure of that response field, and can it tell us which branches to evaluate exactly?}
We answer it on single-track timetabling, building on the exact B\&B of Zhou and Zhong
(2007) and the time-space Lagrangian framework of Meng and Zhou (2014).

\paragraph{Contributions.}
\begin{enumerate}
\item A faithful, validated single-track timetabling B\&B (meet--pass precedence branching;
crossing-conflict enhanced lower bound) used as a controlled data-generating process
(Section~\ref{sec:bb}).
\item A \emph{compressed-response screening} framework with a \emph{gradient-weighted} score
and exact certification of a statistical tie cluster $K'$; we show why the unsupervised
spectral footprint norm is the wrong object (Section~\ref{sec:screen}).
\item A learned latent / low-rank tensor realization of the screen, and an honest empirical
study showing it transfers across instances but that, on single-track, the win is fidelity
and lower-bound node reduction rather than wall-clock (Sections~\ref{sec:learn}--\ref{sec:exp}).
\end{enumerate}

\section{Related work and lineage}
\paragraph{Single-track B\&B (the engine we use).} Zhou and Zhong (2007) formulate
single-track timetabling as a resource-constrained scheduling problem and solve it by B\&B
that branches on the meet/pass order of the trains in conflict at the earliest contested
resource (adding a finish-to-start precedence), bounds each node by a longest-path schedule,
and fathoms with a \emph{crossing-conflict (CC) enhanced lower bound}
$\mathrm{LB}=\sum_i e_i + \sum_{i,i'}\Theta(i,i')$, $\Theta\in\{h,\,g{+}h,\,2h\}$, derived
from the dwell windows of meeting trains. We re-implement this engine faithfully.

\paragraph{$N$-track time-space Lagrangian relaxation (the solution-method lineage).}
Meng and Zhou (2014) reformulate simultaneous routing and rescheduling on an $N$-track
network with \emph{cumulative-flow} variables $y_f(i,j,t)$ and an additive cell-capacity
constraint $\sum_f y_f(i,j,t)\le \mathrm{Cap}(i,j,t)$. Their solution \emph{dualizes the cell
capacity} with prices $\rho_{l,t}$, which \emph{decouples} the relaxed problem into per-train
\emph{time-dependent least-cost path} subproblems on an extended time-space network; the
multipliers are updated by \emph{subgradient} steps, a feasible upper bound is recovered by a
\emph{priority-rule} heuristic (trains ordered by Lagrangian profit), and relax-and-cut
generates capacity constraints dynamically. This additive cell-capacity structure is exactly
the resource-time incidence ($H$ matrix) our screen operates on, and the dual prices
$\rho_{l,t}$ are the resource sensitivities that the gradient weighting requires.

\paragraph{Learning and spectral methods.} Strong, pseudo-cost, and reliability branching
trade accuracy for cost; learning-to-branch (e.g.\ GNN imitation of strong branching) ranks
candidates from data but does not preserve certificates. Low-rank/spectral methods summarize
large incidence matrices. Our framework keeps exact certification and makes the screening
object the \emph{gradient-weighted} low-rank response of the branching action, not a generic
matrix factorization.

\section{Single-track model and the branch-and-bound engine}\label{sec:bb}
\paragraph{Instance.} A single-track line of $S$ sections and two directions; each train has a
type (fast/slow) fixing its per-section running time, a release time, and a direction. A
single headway constant $g_I$ governs both segment and station separation. Resource-time
cells are (section,\,bin) pairs; a train occupies a cell during its traversal plus a headway
tail. Capacity is one train per section-bin, so two opposing (or overtaking) trains cannot
overlap on a section---the meet--pass constraint.

\paragraph{State and schedule generation.} A node holds the timetable
$\textsc{Tab}[\text{train}][\text{section}]=(\text{arrive},\text{depart},\text{floor},\text{fix})$;
``waiting'' is encoded by a raised entry floor, i.e.\ departure differences (the
timetable/Way-2 representation). A forward sweep computes arrival/departure from start times,
running times, and the floors.

\paragraph{Branching.} The earliest-completing active task's section, if contended by
$\ge 2$ trains within the headway window, induces a meet--pass branch: one child per
contending train (``that train goes first''); losers' floors are pushed past the winner's
clearance$+g_I$ and re-swept.

\paragraph{Lower bound.} The CC bound counts only \emph{unresolved} crossing conflicts on
unfixed sections, each contributing $\ge g_I$ over disjoint train pairs---a valid bound (we
verified that the naive $g_I\times(\text{all conflicts})$ form over-prunes; restricting to
unfixed disjoint conflicts restores validity).

\paragraph{Validation.} The CC-bounded search returns the same optimum as the unbounded
search on every tested instance ($6$--$14$ trains, multiple seeds), while reducing nodes
(Table~\ref{tab:lb}). This confirms the engine reproduces the intended exact behavior.

\begin{table}[t]\centering
\caption{Crossing-conflict lower bound: B\&B node reduction vs.\ the unbounded search
(mean over $5$ seeds, exact optimum preserved). Node reduction grows with problem size; at
$16$ trains the unbounded search is intractable (node cap hit), i.e.\ the bound matters most
exactly where the naive search blows up.}\label{tab:lb}
\begin{tabular}{rrrr}
\toprule
trains & nodes (no LB) & nodes (CC-LB) & node reduction \\
\midrule
6  & 108     & 96     & 7.5\% \\
8  & 1{,}053 & 874    & 15.1\% \\
10 & 11{,}900& 8{,}588& 24.7\% \\
12 & 139{,}169 & 91{,}929 & 30.5\% \\
14 & 116{,}910 & 75{,}296 & 35.6\% \\
16 & \multicolumn{3}{c}{intractable unbounded (node cap)} \\
\bottomrule
\end{tabular}
\end{table}

\section{Compressed-response screening for branching}\label{sec:screen}
At a node, let $\mathcal{A}$ be the candidate meet--pass actions. Each action $a$ produces a
child whose re-optimized schedule is the \emph{response} $x(a)\in\mathbb{R}^L$ and whose bound
is the \emph{score} $S(x(a))$. Full strong branching evaluates $S$ for all $a$; we screen
cheaply and certify only a few.

\paragraph{Low-rank response field.} Stacking responses $V=[x(a_1),\dots,x(a_m)]$ and taking
the SVD $V_c=U\Sigma W^\top$, the squared singular values $\sigma_i^2$ are response variances;
a rank-$r$ proxy reconstructs any response in $r$ latent coordinates with error confined to
the tail $\sum_{i>r}\sigma_i^2$.

\paragraph{Gradient-weighted score error (the key result).} We rank by $S$, not by $x$, so the
relevant error is in the score. A first-order expansion gives
$S(x(a))-S(\hat x(a))\approx \nabla S^\top \delta(a)$, and for the total-completion objective
$S=c^\top x$ this is \emph{exact}. Hence the score noise is the \emph{gradient-weighted tail}
\begin{equation}
\mathrm{Var}\!\left(S\text{-err}\right)\;\approx\;\sum_{i>r}\sigma_i^2\,(u_i^\top\nabla S)^2
\;\le\; \|\nabla S\|^2\sum_{i>r}\sigma_i^2 .
\label{eq:gw}
\end{equation}
A discarded direction hurts the ranking \emph{only if the objective gradient points along it}.
Consequently the screening score is not the footprint norm $\|\Sigma_k V_k^\top b_a\|$ but the
gradient-weighted combination
\begin{equation}
\hat S(a)\;\approx\;\sum_i \sigma_i\,(u_i^\top \nabla S)\,(v_i^\top b_a),
\label{eq:score}
\end{equation}
which in practice is realized by a learned score that approximates $S$ rather than response
magnitude.

\paragraph{The statistical tie cluster.} Modeling proxy scores as
$\hat S_a=S_a+\varepsilon_a,\ \varepsilon_a\sim\mathcal N(0,\sigma_{\text{score}}^2(r))$ with
$\sigma_{\text{score}}(r)=\sqrt{\sum_{i>r}\sigma_i^2(u_i^\top\nabla S)^2}$, a worse action can
outrank the best only within the proxy noise, so the \emph{kept set} is the tie cluster
\begin{equation}
\mathcal A_K=\{\,a:\ \hat S_a-\hat S_{\text{best}}\le 2\,\sigma_{\text{score}}(r)\,\},\qquad K'=|\mathcal A_K|.
\end{equation}
Only $\mathcal A_K$ is certified by exact child evaluation; the rank is chosen by the inverse
rule $r^\star=\min\{r:\sigma_{\text{score}}(r)\le\tau\}$.

\begin{proposition}[Screening validity]
If every action in $\mathcal A_K$ is evaluated by the same exact child procedure as full strong
branching, then each evaluated child bound is valid, and the selected action equals the
full-strong-branching action whenever the true best lies in $\mathcal A_K$. When
$\mathcal A_K=\mathcal A$ the method reduces to full strong branching. The spectral/learned layer
changes only which candidates are inspected, never the child relaxation or bound validity.
\end{proposition}

\paragraph{Evidence for \eqref{eq:gw} on train scheduling.} On a root node ($6$ trains), the
response variances are $[2265,1083,264,86,40]$ (95\% energy at rank $3$), yet the
gradient-weighted score noise only collapses at rank $4$ ($\sigma_{\text{score}}:9.37\to0.62$,
realized RMSE $3.82\to0.25$): a rank-$4$ direction with only $2.3\%$ of the energy carries the
score-relevant variance. The scalar bound $\|\nabla S\|\sqrt{\mathrm{tail}}$ overstates the
realized error by $5$--$20\times$. At the energy rank the $2\sigma$ band exceeds all score gaps,
so $K'=$ all moves; raising the rank shrinks $K'$ to $\approx 2$. This is precisely
\eqref{eq:gw}: energy rank $\neq$ score rank.

\section{Learned and tensor realization}\label{sec:learn}
Across instances we collect, per node and candidate action, cheap parent-side features in
layers---node (depth/gap), action (type, overlap, terminal flag), dual proxies (contention,
occupancy, slack), trajectory (train completion/delay/conflict count)---and a multi-metric
response $Y=[\Delta\mathrm{LB}_L,\Delta\mathrm{LB}_R,\Delta\mathrm{Strong},\Delta\mathrm{UB},
\text{imbalance},\Delta\mathrm{Primal},\Delta\mathrm{Violation}]$, with a per-section response
tensor. A learned score $\hat S(a)=f_\theta(\cdot)$ realizes the gradient-weighted screen;
reduced-rank multi-output regression realizes the tensor (CP/Tucker) view. This is \emph{not}
$H/B\to\mathrm{SVD}\to\text{label}$: the incidence skeleton is lifted into a multi-layer
feature tensor before any decomposition.

\section{Computational results (single-track)}\label{sec:exp}
\paragraph{Node-level screening.} On $1{,}279$ test nodes, ranking each method's top-$K$ against
the full-strong-branching best action (Table~\ref{tab:contain}), the learned gradient-aware
score is best at every $K$ and the unsupervised spectral norm is worst---direct confirmation
that screening must be gradient-weighted, not energy-weighted. Containment@$K$ is the empirical
tie cluster $K'$: the best branch is rarely top-$1$ but almost always within the top $\approx 5$.

\begin{table}[t]\centering
\caption{Node-level screening vs.\ full strong branching ($1{,}279$ nodes).}\label{tab:contain}
\begin{tabular}{lcccc}
\toprule
method & Contain@1 & Contain@3 & Contain@5 & mean RankGap \\
\midrule
learned (gradient-aware) & \bf 0.37 & \bf 0.77 & \bf 0.92 & \bf 1.56 \\
unsupervised spectral norm & 0.17 & 0.56 & 0.77 & 2.80 \\
pseudo-cost (overlap) & 0.27 & 0.63 & 0.80 & 2.40 \\
random & 0.27 & 0.67 & 0.85 & 2.15 \\
\bottomrule
\end{tabular}
\end{table}

\paragraph{Transfer across instances.} Trained on $14$ instances and tested on unseen
instances (Table~\ref{tab:transfer}), the screen generalizes (out-of-instance Containment@5
$\approx 0.5$--$0.63$, Spearman $\approx 0.8$), far above pseudo-cost/random ($\approx 0$). A
twist: the \emph{dual} layer massively raises global rank-correlation (Spearman $+0.1\to+0.5$)
but does \emph{not} improve per-node Containment@5---global ordering $\neq$ local top-$K$. The
\emph{trajectory} (train-delay) layer is the consistent winner on both metrics; permutation
importance attributes $44\%$--$60\%$ to the trajectory layer and only $1$--$2\%$ to the dual
proxies (single-track).

\begin{table}[t]\centering
\caption{Out-of-instance transfer (size $10$): feature subsets predicting $\Delta\mathrm{Strong}$.}
\label{tab:transfer}
\begin{tabular}{lcc}
\toprule
feature set & Containment@5 & Spearman (unseen) \\
\midrule
action-only & 0.35 & +0.11 \\
action + dual & 0.34 & +0.52 \\
action + trajectory & 0.51 & +0.78 \\
all layers & \bf 0.63 & \bf +0.85 \\
pseudo-cost / random & 0.33 / 0.27 & $\approx 0$ \\
\bottomrule
\end{tabular}
\end{table}

\paragraph{Tensor structure.} The per-section response tensor is low multilinear rank
(metric mode \emph{rank one}, $\approx 98\%$ energy; HOSVD rank $(5,4,1)/(6,5,1)$); it is
essentially a single interior wait-time field rather than a clean taxonomy of congestion
modes. The coupled $7$-metric response is low-rank (rank $2$ of $7$ on the larger size), giving
\emph{parsimony}; but low-rank multi-output regression does not beat a single-output nonlinear
predictor of $\Delta\mathrm{Strong}$---the tensor view earns interpretability and transfer, not
accuracy.

\paragraph{Honest scope.} On single-track the candidate set is small (mean $|\mathcal A|\approx
5.4$) and the cheap earliest-conflict rule already minimizes the tree; strong-branching-style
selection (and its learned proxy) does \emph{not} reduce node count. Screening the tie cluster
cuts exact child evaluations only $1.3$--$1.9\times$. The durable single-track gains are
(i) the CC lower bound's $7.5$--$35.6\%$ node reduction (Table~\ref{tab:lb}) and (ii) screening
\emph{fidelity}: contain the optimal branch in a small, exactly-certified tie cluster.

\section{Outlook: the $N$-track regime}
The single-track study isolates the mechanism but caps the payoff: tie clusters are large,
candidate sets small, and the dual proxies are weak surrogates for true prices. The speedup
regime is the $N$-track space-time network of Meng and Zhou (2014), where route/track choice
enlarges $\mathcal A$ and the LP relaxation supplies \emph{true} cell-time duals
$\rho_{l,t}$---the exact quantity \eqref{eq:gw} needs. We have a working space-time LP on the
INFORMS RAS network ($85$ arcs, $20$ trains, $\sim$11{,}400 trajectory columns, $\sim$197{,}900
cell-time rows; $142$ fractional variables and $128$ binding cell-time duals at the root),
which is the certified next step for testing whether gradient-weighted screening yields a
genuine branch-and-price speedup.

\section{Conclusion}
Treating the timetabling solver as a data-generating process, we showed that the low-rank
structure of its response field can screen branching candidates, provided the screen is
\emph{gradient-weighted} and only a statistical tie cluster is certified exactly. On
single-track the mechanism is real and transfers across instances, the unsupervised spectral
score provably underperforms, and the crossing-conflict lower bound delivers a $7.5$--$35.6\%$
node reduction; a per-solve speedup awaits the larger-candidate, true-dual $N$-track regime.

\begin{thebibliography}{9}
\bibitem{ZhouZhong2007} X.~Zhou and M.~Zhong. Single-track train timetabling with guaranteed
optimality: Branch-and-bound algorithms with enhanced lower bounds. \emph{Transportation
Research Part B}, 41(3):320--341, 2007.
\bibitem{MengZhou2014} L.~Meng and X.~Zhou. Simultaneous train rerouting and rescheduling on an
$N$-track network: A model reformulation with network-based cumulative flow variables.
\emph{Transportation Research Part B}, 67:208--234, 2014.
\bibitem{Caprara2002} A.~Caprara, M.~Fischetti, and P.~Toth. Modeling and solving the train
timetabling problem. \emph{Operations Research}, 50(5):851--861, 2002.
\bibitem{Achterberg2005} T.~Achterberg, T.~Koch, and A.~Martin. Branching rules revisited.
\emph{Operations Research Letters}, 33(1):42--54, 2005.
\bibitem{Khalil2016} E.~Khalil et al. Learning to branch in mixed integer programming.
\emph{AAAI}, 2016.
\bibitem{Gasse2019} M.~Gasse et al. Exact combinatorial optimization with graph convolutional
neural networks. \emph{NeurIPS}, 2019.
\bibitem{BLP2021} Y.~Bengio, A.~Lodi, and A.~Prouvost. Machine learning for combinatorial
optimization: A methodological tour d'horizon. \emph{EJOR}, 290(2):405--421, 2021.
\end{thebibliography}

\end{document}
